\section{图解：三角换元的几何关系}
\begin{figure}[htbp]\centering
\begin{tikzpicture}[scale=1.2,>=stealth]
\draw[thick,Ink](0,0)--(4,0)--(4,2.4)--cycle;
\draw(3.72,0)--(3.72,.28)--(4,.28);
\draw[Teal,->](.7,0) arc[start angle=0,end angle=30.96,radius=.7];
\node at (1,.24){$\theta$};
\node[below]at(2,0){$a>0$};\node[right]at(4,1.2){$x$};
\node[above,rotate=31]at(2,1.2){$\sqrt{a^2+x^2}$};
\node[align=left,anchor=west]at(5,.9){$\tan\theta=\frac{x}{a}$\\$x=a\tan\theta$\\$\dd x=a\sec^2\theta\,\dd\theta$};
\end{tikzpicture}
\caption{对$x>0$，直角三角形给$\sqrt{a^2+x^2}=a\sec\theta$。对负$x$用$\theta\in(-\frac{\pi}{2},0)$的有向变量仍有$\sec\theta>0$，故全体实$x$可用$\theta=\arctan(\frac{x}{a})$。图形只帮助识别恒等式，微分因子仍须实际计算。}\label{fig:review:substitution}
\end{figure}
